% Shared layout for both CV versions (based on Jake Gutierrez's resume template, MIT).
\documentclass[letterpaper,11pt]{article}

\usepackage[utf8]{inputenc}
\usepackage[T1]{fontenc}
\usepackage{latexsym}
\usepackage[empty]{fullpage}
\usepackage{titlesec}
\usepackage[usenames,dvipsnames]{color}
\usepackage{enumitem}
\usepackage[hidelinks]{hyperref}
\usepackage{fancyhdr}
\usepackage{tabularx}
\usepackage{fontawesome5}
\input{glyphtounicode}

\pagestyle{fancy}
\fancyhf{}
\fancyfoot{}
\renewcommand{\headrulewidth}{0pt}
\renewcommand{\footrulewidth}{0pt}

\addtolength{\oddsidemargin}{-0.5in}
\addtolength{\evensidemargin}{-0.5in}
\addtolength{\textwidth}{1in}
\addtolength{\topmargin}{-.5in}
\addtolength{\textheight}{1.0in}

\urlstyle{same}
\raggedbottom
\raggedright
\setlength{\tabcolsep}{0in}

\titleformat{\section}{\vspace{-4pt}\scshape\raggedright\large\bfseries}{}{0em}{}[\color{black}\titlerule \vspace{-5pt}]

\pdfgentounicode=1

\newcommand{\resumeItem}[1]{\item\small{{#1 \vspace{-2pt}}}}

\newcommand{\resumeSubheading}[4]{
  \vspace{-2pt}\item
    \begin{tabular*}{0.97\textwidth}[t]{l@{\extracolsep{\fill}}r}
      \textbf{#1} & #2 \\
      \textit{\small#3} & \textit{\small #4} \\
    \end{tabular*}\vspace{-7pt}
}

\newcommand{\resumeProjectHeading}[2]{
    \item
    \begin{tabular*}{0.97\textwidth}{p{0.75\textwidth}@{\extracolsep{\fill}}r}
      \small#1 & \small #2 \\
    \end{tabular*}\vspace{-7pt}
}

\newcommand{\resumeAward}[3]{
    \item
    \begin{tabular*}{0.97\textwidth}{l@{\extracolsep{\fill}}r}
      \small\textbf{#1} & \small #2 \hspace{1.5em} #3 \\
    \end{tabular*}\vspace{-7pt}
}

\renewcommand\labelitemii{$\vcenter{\hbox{\tiny$\bullet$}}$}
\newcommand{\resumeSubHeadingListStart}{\begin{itemize}[leftmargin=0.15in, label={}]}
\newcommand{\resumeSubHeadingListEnd}{\end{itemize}}
\newcommand{\resumeItemListStart}{\begin{itemize}}
\newcommand{\resumeItemListEnd}{\end{itemize}\vspace{-5pt}}

\newcommand{\cvheader}[1]{
\begin{center}
    {\Huge \scshape Juan del Mazo Vázquez} \\ \vspace{3pt}
    \small \faPhone*\ \href{tel:+34678641703}{(+34) 678 64 17 03} \quad
    \faEnvelope\ \href{mailto:delmazo.vazquez@gmail.com}{delmazo.vazquez@gmail.com} \quad
    \faLinkedin\ \href{https://www.linkedin.com/in/juan-del-mazo-7500a7276}{juan-del-mazo} \quad
    \faGithub\ \href{https://github.com/JuandelMazoVazquez}{JuandelMazoVazquez} \\ \vspace{2pt}
    \faGlobe\ \href{https://juandelmazovazquez.vercel.app}{juandelmazovazquez.vercel.app} \quad
    \faMapMarker*\ #1
\end{center}
}


\begin{document}

\cvheader{Madrid, España}
  {Analista de profesión, ingeniero de formación.}
  {Sociable, proactivo y curioso. Me gusta construir cosas, decidir con datos y dar un paso al frente cuando toca liderar.}

\section{Experiencia profesional}
  \resumeSubHeadingListStart
    \resumeSubheading
      {Junior Business Analyst}{sept. 2025 -- actualidad}
      {Bankinter -- Análisis de Préstamos Preautorizados de Consumo}{Madrid}
      \resumeItemListStart
        \resumeItem{Analizo datos para apoyar las decisiones sobre la cartera de préstamos preautorizados de consumo, con un doble objetivo: maximizar la producción y minimizar el riesgo.}
        \resumeItem{Hago análisis de segmentación y seguimiento del riesgo para detectar margen de crecimiento y vigilar cómo evoluciona el riesgo de la cartera.}
        \resumeItem{Preparo propuestas para optimizar el negocio y las presento a los directores de Análisis y de Consumo.}
      \resumeItemListEnd

    \resumeSubheading
      {Desarrollador de software}{jun. 2023 -- jun. 2024}
      {Idaero Solutions S.L.}{Software para el sector aeroespacial}
      \resumeItemListStart
        \resumeItem{Dirigí el desarrollo de la interfaz de visualización de un gemelo digital que predice defectos en la fabricación de materiales compuestos para el sector aeroespacial.}
        \resumeItem{Coordiné a stakeholders de 5 empresas para alinear requisitos técnicos e hitos del proyecto.}
        \resumeItem{Desarrollé en menos de 3 meses una librería de bajo nivel de 8.000+ líneas que usa la aplicación principal.}
        \resumeItem{Corregí errores críticos y menores y di soporte técnico.}
      \resumeItemListEnd

    \resumeSubheading
      {Becario de desarrollo de software}{feb. 2023 -- jun. 2023}
      {Idaero Solutions S.L.}{Software para el sector aeroespacial}
      \resumeItemListStart
        \resumeItem{Desarrollé en C\# funcionalidades de alto impacto (6.000+ líneas) en un equipo ágil, con GitHub para el control de versiones.}
      \resumeItemListEnd
  \resumeSubHeadingListEnd

\section{Formación}
  \resumeSubHeadingListStart
    \resumeSubheading
      {Máster en Big Data Analytics}{2024 -- 2025}
      {Universidad Carlos III de Madrid -- Calificación final: 1st Class (Hons.); impartido en inglés}{Madrid}
    \resumeSubheading
      {Grado en Ingeniería Informática}{2019 -- 2024}
      {Universidad Politécnica de Madrid -- Nota media: 1st Class (Hons.)}{Madrid}
    \resumeSubheading
      {Erasmus en Informática (Computer Science)}{2021 -- 2022}
      {Universidad de Bristol -- Nota media: 1st Class (Hons.)}{Bristol, Reino Unido}
  \resumeSubHeadingListEnd
  \vspace{-2pt}
  \begin{itemize}[leftmargin=0.15in, label={}]
    \small\item \textbf{Matrícula de Honor} en Web Data Analytics and Usage (UC3M, 2025), Procesadores de Lenguajes (UPM, 2022), Matemática Discreta II, Arquitectura de Computadores y Estructura de la Concurrencia (UPM, 2021).
  \end{itemize}

\section{Proyectos}
  \resumeSubHeadingListStart
    \resumeProjectHeading{\textbf{TFM: efectos lead-lag en trading} $|$ \emph{Python, DTW, clustering espectral, backtesting} \href{https://github.com/JuandelMazoVazquez/Lead-Lag-Trading-Strategy}{\faGithub}}{ene. 2025 -- jul. 2025}
      \resumeItemListStart
        \resumeItem{Construí un sistema no supervisado que detecta qué acciones se adelantan a las demás entre 500 valores de EE.\,UU. en 20 años (4.700+ ventanas móviles) y lo convertí en una señal de trading: ratio de Sharpe de 0,596, un 35\,\% por encima del modelo de referencia (0,442).}
      \resumeItemListEnd
    \resumeProjectHeading{\textbf{Shuffl: modo aleatorio de verdad para Spotify} $|$ \emph{JavaScript, PWA, Playwright} \href{https://github.com/JuandelMazoVazquez/new-spoti}{\faGithub}}{2026}
      \resumeItemListStart
        \resumeItem{Diseñé, construí y publiqué una aplicación web sin servidor con un Fisher--Yates sin sesgo; va por la v1.4 y tiene 300+ pruebas unitarias y de extremo a extremo en CI.}
      \resumeItemListEnd
    \resumeProjectHeading{\textbf{Canal de vídeo automatizado (en marcha)} $|$ \emph{Python, Gemini, Azure Speech, FFmpeg}}{2026}
      \resumeItemListStart
        \resumeItem{Sistema con IA que convierte hechos verificados en vídeos cortos (guion escrito por un LLM, síntesis de voz y montaje automático) y los publica dos veces al día en YouTube y TikTok; 20.000+ líneas de Python y 6.900+ pruebas.}
      \resumeItemListEnd
    \resumeProjectHeading{\textbf{Casos prácticos del máster (en equipo)} $|$ \emph{R, Python, XGBoost, SHAP, Pyomo}}{2024 -- 2025}
      \resumeItemListStart
        \resumeItem{Previsión de ventas según la inversión en marketing (ensemble ponderado, MAE de validación de 298,75), rotación de empleados con datos desbalanceados (balanced accuracy de 0,92) y optimización estocástica en dos etapas con hasta 301 escenarios.}
      \resumeItemListEnd
  \resumeSubHeadingListEnd

\section{Competencias}
 \begin{itemize}[leftmargin=0.15in, label={}]
    \small{\item{
     \textbf{Dominio alto}{: Python, C\#, SQL, Java} \quad \textbf{Avanzado}{: R, Spark, C++, C} \quad \textbf{Intermedio}{: MATLAB, Bash, ensamblador, JavaScript, TypeScript} \\
     \textbf{Librerías y herramientas}{: pandas, scikit-learn, XGBoost, PySpark, Pyomo, Playwright, Git, GitHub Actions} \\
     \textbf{IA}{: APIs de LLM (Gemini), Claude Code, T3 Code, flujos multiagente, prompt engineering, síntesis de voz, automatización con IA} \\
     \textbf{Áreas}{: análisis de datos, riesgo de crédito, aprendizaje automático, finanzas cuantitativas, optimización, desarrollo de software} \\
     \textbf{Idiomas}{: español (nativo), inglés (nivel alto)}
    }}
 \end{itemize}

\section{Otros}
 \begin{itemize}[leftmargin=0.15in, label={}]
    \small{\item{
     Clases particulares de Matemáticas (ESO y Bachillerato) y de Arquitectura de Computadores a universitarios (2021 -- 2023) $\cdot$
     Becario en la Delegación de Alumnos de la UPM: llevé los créditos de 92 estudiantes (2022) $\cdot$
     Organicé y dirigí un torneo de pádel de 26 jugadores (2022).
    }}
 \end{itemize}

\end{document}
